\subsection{Divisorial ideals of an integral domain}

DEFINITION 1. Let A be an integral domain and K its field of fractions. Every sub-A-module a of K such that there exists an element \(d \neq 0\) in A for which da ⊂ A is called a fractional ideal of A (or of K, by an abuse of language).

Every finitely generated sub-A-module \(\mathfrak{a}\) of \(\mathbf{K}\) is a fractional ideal: for if \((a_{i})_{1\leqslant i\leqslant n}\) is a system of generators of \(\mathfrak{a}\) , we may write \(a_{i} = b_{i} / d_{i}\) , where \(b_{i}\in \mathbf{A}\) , \(d_{i}\in \mathbf{A}\) and \(d_{i}\neq 0\) ; if \(d = d_{1}\ldots d_{n}\) , clearly \(da\subset A\) . In particular the monogenous sub-A-modules of \(\mathbf{K}\) are fractional ideals (recall that they have been called fractional principal ideals in Algebra, Chapter VI, § 1, no. 5). If \(\mathbf{A}\) is Noetherian, every fractional ideal is a finitely generated A-module. Every sub-A-module of a fractional ideal of \(\mathbf{A}\) is a fractional ideal. Every ideal of \(\mathbf{A}\) is a fractional ideal; to avoid confusion, these will also be called the integral ideals of \(\mathbf{A}\) .

We denote by I(A) the set of non-zero fractional ideals of A. Given two elements a, b of I(A), we shall write \(a \prec b\) (or \(b \succ a\) ) for the relation ``every fractional principal ideal containing a also contains b''; clearly this relation is a preordering on I(A). Let R denote the associated equivalence relation ``a \(\prec b\) and \(b \prec a\) '' (Set Theory, Chapter III, § 1, no. 2) and D(A) the quotient set I(A)/R; we shall say that the elements of D(A) are the divisors of A and, for every fractional ideal \(a \in I(A)\) , we shall denote by div a (or div \(_{A}\) a) the canonical image of a in D(A) and we shall say that div a is the divisor of a; if a = Ax is a fractional principal ideal, we write \(\operatorname{div}(x)\) instead of \(\operatorname{div}(\mathbf{A}x)\) and \(\operatorname{div}(x)\) is called the divisor of \(x\) ; the elements of D(A) of the form \(\operatorname{div}(x)\) are called principal divisors. By taking the quotient, the preordering \(\prec\) on I(A) defines on D(A) an ordering which we shall denote by \(\leqslant\) .

For all \(\mathfrak{a} \in \mathrm{I}(\mathbf{A})\) there exists by hypothesis some \(d \neq 0\) in \(\mathbf{A}\) such that \(\mathfrak{a} \subset \mathrm{Ad}^{-1}\) ; the intersection \(\tilde{\mathfrak{a}}\) of the fractional principal ideals containing \(\mathfrak{a}\) is therefore an element of \(\mathrm{I}(\mathbf{A})\) . Clearly the relation \(\mathfrak{a} \prec \mathfrak{b}\) is equivalent to the relation \(\tilde{\mathfrak{a}} \supset \tilde{\mathfrak{b}}\) ; the relation \(\mathfrak{a} \supset \mathfrak{b}\) therefore implies \(\mathfrak{a} \prec \mathfrak{b}\) . For two elements \(\mathfrak{a}, \mathfrak{b}\) of \(\mathrm{I}(\mathbf{A})\) to be equivalent modulo \(\mathbb{R}\) , it is necessary and sufficient that \(\tilde{\mathfrak{a}} = \tilde{\mathfrak{b}}\) .

DEFINITION 2. Every element \(\mathfrak{a}\) of I(A) such that \(\mathfrak{a} = \tilde{\mathfrak{a}}\) is called a divisorial fractional ideal of A.

In other words a divisorial ideal is just a non-zero intersection of a non-empty family of fractional principal ideals. Every non-zero intersection of divisorial ideals is a divisorial ideal. If a is divisorial, so is ax for all \(x \in K^{*}\) , the mapping \(b \mapsto bx\) being a bijection of the set of fractional principal ideals onto itself. For all \(a \in I(A)\) , \(\tilde{a}\) is the least divisorial ideal containing a and is equivalent to a modulo R. Moreover, if b is a divisorial ideal equivalent to a modulo R, then \(\tilde{a} = \tilde{b} = b\) . Hence \(\tilde{a}\) is the unique divisorial ideal b such that div a = div b (in other words, the restriction of the mapping \(a \mapsto div a\) to the set of divisorial ideals is injective).

Let \(\mathfrak{a}\) and \(\mathfrak{b}\) be two fractional ideals of \(\mathbf{K}\) . Recall (Chapter I, § 2, no. 10) that \(\mathfrak{b}\) : \(\mathfrak{a}\) denotes the set of \(x \in \mathbf{K}\) such that \(x\mathfrak{a} \subset \mathfrak{b}\) ; this is obviously an A-module; if \(\mathfrak{b} \in \mathrm{I}(\mathrm{A})\) and \(\mathfrak{a} \in \mathrm{I}(\mathrm{A})\) , then \(\mathfrak{b} \colon \mathfrak{a} \in \mathrm{I}(\mathrm{A})\) ; for if \(d\) is a non-zero element of \(\mathbf{A}\) such that \(db \subset \mathbf{A}\) and \(da \subset \mathbf{A}\) and \(a\) is a non-zero element of \(\mathbf{A} \cap \mathfrak{a}\) , then \(da(\mathfrak{b} : \mathfrak{a}) \subset \mathbf{A}\) ; on the other hand, if \(b \neq 0\) belongs to \(\mathfrak{b}\) , then \(bda \subset \mathfrak{b}\) , hence \(bd \in \mathfrak{b}\) : \(\mathfrak{a}\) and \(\mathfrak{b} : \mathfrak{a} \neq 0\) .

The definition of b: a can also be written:

\[
\mathfrak {b}: \mathfrak {a} = \bigcap_ {x \in \mathfrak {a}, x \neq 0} \mathfrak {b} x ^ {- 1}.\tag{1}
\]

PROPOSITION 1. (a) If \(\mathfrak{b}\) is a divisorial ideal and \(\mathfrak{a} \in \mathbf{I}(\mathbf{A})\) , \(\mathfrak{b}\) : \(\mathfrak{a}\) is divisorial.

\begin{enumerate}
\def\labelenumi{(\alph{enumi})}
\setcounter{enumi}{1}
\item
  Let \(\mathfrak{a},\mathfrak{b}\) be in I(A). In order that \(\operatorname {div}\mathfrak{a} = \operatorname {div}\mathfrak{b}\) , it is necessary and sufficient that A: \(\mathfrak{a} = \mathbf{A}\) : \(\mathfrak{b}\) .
\item
  For all \(\mathfrak{a} \in \mathbf{I}(\mathbf{A})\) , \(\tilde{\mathfrak{a}} = \mathbf{A} : (\mathbf{A} : \mathfrak{a})\) .
\end{enumerate}

Assertion (a) follows immediately from equation (1) since, if \(\mathfrak{b}\) is divisorial, so is \(\mathfrak{bx}^{-1}\) for all \(x \neq 0\) .

To show (b), let \(\mathrm{P}(\mathfrak{a})\) denote the set of fractional principal ideals containing \(\mathfrak{a}\) ; the relation \(Ax \in \mathrm{P}(\mathfrak{a})\) is equivalent to \(x^{-1}\mathfrak{a} \subset A\) and hence to \(x^{-1} \in A\) : \(\mathfrak{a}\) . As the relation \(\operatorname{div} \mathfrak{a} = \operatorname{div} \mathfrak{b}\) is by definition equivalent to \(\mathrm{P}(\mathfrak{a}) = \mathrm{P}(\mathfrak{b})\) , it is also equivalent to \(A\) : \(\mathfrak{a} = A\) : \(\mathfrak{b}\) .

Finally, as \(\mathfrak{a}(\mathbf{A}:\mathfrak{a})\subset \mathbf{A},\mathfrak{a}\subset \mathbf{A}:(\mathbf{A}:\mathfrak{a})\) . Replacing \(\mathfrak{a}\) by \(\mathbf{A}:\mathfrak{a}\) in this formula, it is seen that \(\mathbf{A}:\mathfrak{a}\subset \mathbf{A}:(\mathbf{A}:(\mathbf{A}:\mathfrak{a}))\) ; on the other hand, the relation \(\mathfrak{a}\subset \mathbf{A}:(\mathbf{A}:\mathfrak{a})\) implies

\[
\mathrm{A}: \mathfrak {a} \supset \mathrm{A}: (\mathrm{A}: (\mathrm{A}: \mathfrak {a})).
\]

Therefore A: a = A : (A : (A : a)) and it follows from (b) that div a = div(A : (A : a)) As A : (A : a) is divisorial by (a), certainly \(\tilde{a} = A\) : (A : a), which proves (c).

Remark. In the course of the above proof it has been proved that A: a = A: (A: (A: a)) for every ideal \(a \in I(A)\) , which is a special case of Set Theory, Chapter III, § 1, no. 5, Proposition 2.

PROPOSITION 2. (i) In D(A) every non-empty set bounded above admits a least upper bound. More precisely, if \((\mathfrak{a}_t)\) is a non-empty family of elements of I(A) which is bounded above, then

\[
\sup _ {i} (\operatorname{div} a _ {i}) = \operatorname{div} \left(\bigcap_ {i} \tilde {a} _ {i}\right).
\]

\begin{enumerate}
\def\labelenumi{(\roman{enumi})}
\setcounter{enumi}{1}
\tightlist
\item
  In D(A) every non-empty set bounded below admits a greatest lower bound. More precisely, if \((\mathfrak{a}_{\mathfrak{t}})\) is a non-empty family of elements of I(A) which is bounded below, then
\end{enumerate}

\[
\inf _ {i} (\operatorname{div} a _ {i}) = \operatorname{div} \left(\sum_ {i} a _ {i}\right).
\]

\begin{enumerate}
\def\labelenumi{(\roman{enumi})}
\setcounter{enumi}{2}
\tightlist
\item
  The set D(A) is a lattice.
\end{enumerate}

Let \((\mathfrak{a}_i)\) be a non-empty family of elements of I(A) which is bounded above. To say that a divisorial ideal \(\mathfrak{b}\) bounds this family above amounts to saying that it is contained in all the \(\tilde{\mathfrak{a}}_i\) , that is that \(\mathfrak{b}\) is contained in \(\bigcap_{i} \tilde{\mathfrak{a}}_i\) . Hence \(\bigcap_{i} \tilde{\mathfrak{a}}_i \neq (0)\) and \(\bigcap_{i} \tilde{\mathfrak{a}}_i\) is therefore a divisorial ideal, which shows (i).

Now let \((\tilde{a}_t)\) be a non-empty family of elements of I(A) which is bounded below. To say that a divisorial ideal \(b\) bounds this family below means that it contains all the \(\tilde{a}_t\) , that is (since \(b\) is divisorial) that it contains all the \(a_t\) , or also that \(b \supset \sum_{i} a_i\) . This proves (ii).

Finally, to prove (iii) it is sufficient by (i) and (ii) to prove that, if a, b are in I(A), the set \(\{a, b\}\) is bounded both above and below in I(A); now it is bounded above by \(a \cap b\) (which is distinct from (0)). It is bounded below by \(a + b\) , for \(a + b \in I(A)\) : if d and \(d'\) are non-zero elements of A such that \(da \subset A\) and \(d'b \subset A\) , then \(dd'(\mathfrak{a} + \mathfrak{b}) \subset \mathbf{A}\) .

COROLLARY. If \(x, y\) and \(x + y\) are in \(\mathbf{K}^*\) , then \(\operatorname{div}(x + y) \geqslant \inf(\operatorname{div}(x), \operatorname{div}(y))\) . \(\mathbf{A}(x + y) \subset \mathbf{A}x + \mathbf{A}y\) and hence \(\operatorname{div}(x + y) \geqslant \operatorname{div}(\mathbf{A}x + \mathbf{A}y)\) .
% semantic-fragments: legacy-input-seam
\relax{}
